\clearpage
\begin{trapbox}
\textbf{易错、边界与最小反例（C086 专项核查）。}
\begin{itemize}
\item \textbf{把“存在”误读为“对任意点”。}结论是\emph{存在某个点}满足指定比值，
并非双曲线上所有点焦距比相同。例如 $e=2$ 时，焦距比在右支会从 $3$ 连续降向 $1$。

\item \textbf{忽略 $\lambda=1$。}当 $\lambda=1$ 时 $PF_1=PF_2$，但
$|PF_1-PF_2|=2a>0$。\textbf{没有解}，不是上界“无穷大”的特殊情形。

\item \textbf{把 $e$ 的下界写成 $e\ge1$。}非退化双曲线 $b>0$ 导致
$c>a$，因此\textbf{$e>1$，等号 $e=1$ 不可取}。

\item \textbf{将上界的等号漏掉。} $P$ 在对应支顶点时
$PF_1+PF_2=2c$，可以取到\textbf{上界}。
例如 $\lambda=2,e=3$、$P=(a,0)$ 确实满足比值。

\item \textbf{将上下界对调。}正确的是 $1<e\le(1+\lambda)/|1-\lambda|$；
因为\textbf{三角形两边之和 $\ge$ 第三边}。

\item \textbf{漏写绝对值。} $0<\lambda<1$ 时，
$\lambda-1<0$；必须写 $|\lambda-1|$，或分类写分母 $1-\lambda$。

\item \textbf{不能用固定支的焦半径公式套另一支。}
右支：$PF_1=ex+a$、$PF_2=ex-a$；
左支用 $x=-t\le-a$：$PF_1=et-a$、$PF_2=et+a$。

\item \textbf{把焦距写成 $c$。}两个焦点的\textbf{距离为 $2c$}，
而双曲线的焦半距为 $c$；定义中焦距差是 $2a$。

\item \textbf{误用椭圆的定义或参数关系。}双曲线为
$|PF_1-PF_2|=2a$、$c^2=a^2+b^2$；
椭圆才用距离\textbf{和}等于 $2a$ 及 $c^2=a^2-b^2$。

\item \textbf{只证必要不证充分。}从三角形不等式只能推得\emph{必须}满足的范围；
还要用\textbf{两圆交点条件}或\textbf{解得 $|x|\ge a$}，证明范围内确有曲线点。

\item \textbf{误判交点数。}对固定 $\lambda\ne1$，
严格小于离心率上界时有\textbf{两点}；等号时\textbf{一点}；超过时\textbf{零点}。

\item \textbf{误解焦距比为向量比。}$PF_1/PF_2$ 是\textbf{长度之比}，没有方向和正负；
不能以“向量方向”决定 $\lambda$ 的取值。

\item \textbf{把渐近线附近误当等号位置。}焦距比趋近 $1$ 发生于
$|x|\to\infty$，是\textbf{不取到的极限}；上界等号发生在\textbf{顶点}。

\item \textbf{把阿波罗尼斯圆与两焦点辅助圆混为一谈。}
$PF_1=\lambda PF_2$（$\lambda\ne1$）的\textbf{轨迹本身是一条阿波罗尼斯圆}，
其圆心通常\textbf{不在焦点上}；而\textbf{两焦点辅助圆}的圆心是 $F_1,F_2$，
半径按 $2a$ 和 $\lambda$ 计算；不能拿辅助圆的外切条件当作定比圆与双曲线的相切条件。

\item \textbf{误记阿波罗尼斯圆的圆心方向。}
$\lambda>1$ 表示 $PF_1>PF_2$，圆心在\textbf{右侧} $x_C>c$；
$0<\lambda<1$ 时圆心在\textbf{左侧} $x_C<-c$。
$\lambda=1$ 的平面轨迹是\textbf{垂直平分线 $x=0$}，不是圆；
但它依然\textbf{与双曲线无交点}。

\item \textbf{误判两焦点辅助圆的相切类型。}边界时
$r_1+r_2=2c$ 是\textbf{外切}，不是内切；内切对应圆心距等于半径差 $2a$，
但非退化双曲线总有 $2a<2c$。
\end{itemize}

\textbf{自检清单：}
\[
\boxed{\lambda>0,\ \lambda\ne1}\quad
\boxed{e>1}\quad
\boxed{2a=|PF_1-PF_2|}\quad
\boxed{2c=F_1F_2}.
\]
\end{trapbox}
